\title{QLoRA: Efficient Finetuning of Quantized LLMs}

\begin{abstract}
This paper introduces \textsc{QLoRA}, a resource-efficient finetuning method that significantly reduces memory requirements, enabling the finetuning of models with up to 65B parameters on a single 48GB GPU while maintaining the task performance characteristic of full 16-bit finetuning. \textsc{QLoRA} operates by backpropagating gradients through a fixed 4-bit quantized pretrained language model, directing the learning into Low Rank Adapters (LoRA). Our most advanced model series, dubbed \textbf{Guanaco}, establishes new performance records among publicly available models on the Vicuna benchmark, achieving 99.3\% of ChatGPT’s performance after just one day of finetuning on a single GPU. The main innovations of \textsc{QLoRA} that enable these advancements include: (a) the 4-bit NormalFloat (NF4), a novel data format optimized from an information-theoretic perspective for weights with normal distributions, (b) Double Quantization, which further compresses memory by quantizing quantization constants, and (c) Paged Optimizers to regulate episodic memory surges. Utilizing \textsc{QLoRA}, we finetune over 1,000 different models, presenting an extensive evaluation of instruction-following capabilities and conversational aptitude across eight instruction datasets, several backbone architectures (such as LLaMA and T5), and model sizes previously infeasible for standard finetuning methods (including 33B and 65B models). Our experiments reveal that finetuning with QLoRA using concise, high-quality data enables even comparatively compact models to attain state-of-the-art results, surpassing prior best methods. We also provide a comprehensive comparison of chatbot performance using both human judgment and GPT-4-based assessment, showing that GPT-4 offers an affordable and practical proxy for human evaluation. Notably, our investigation suggests that many existing chatbot benchmarks inadequately reflect true performance distinctions. We also conduct a focused error analysis to highlight situations in which \textbf{Guanaco} falls short of ChatGPT. All models and source code, including CUDA implementations supporting 4-bit training, are openly available.\footnote{\url{https://github.com/artidoro/qlora} and \url{https://github.com/TimDettmers/bitsandbytes}}
\end{abstract}

\section{Introduction}

Adapting large language models (LLMs) through finetuning is an established strategy for enhancing their capabilities \citep{min2021metaicl, wei2021finetuned, ouyang2022training, wang2022super, wang2022self, liu2022few} and for modifying or controlling their outputs \citep{ouyang2022training,askell2021general,bai2022training}. Nevertheless, finetuning such massive models imposes substantial computational and financial costs; for instance, conventional 16-bit finetuning of the LLaMA model with 65B parameters~\citep{touvron2023llama} demands in excess of 780 GB of GPU memory. Although advances in quantization have successfully lowered the memory demands for LLM inference \citep{dettmers2022llm, dettmers2022case, frantar2022gptq, xiao2022smoothquant}, these solutions typically do not extend to the training phase, often failing in practice \citep{wortsman2023stable}. 

In this work, we present the first evidence that one can successfully finetune a model quantized to 4 bits without experiencing any loss in performance. Our approach, \textsc{QLoRA}, employs a novel quantization scheme that preserves high precision when mapping a pretrained model to 4 bits, and introduces a compact set of trainable Low-rank Adapter weights \citep{hu2021lora} 

that are optimized by propagating gradients through the quantized model parameters.

By adopting \textsc{QLoRA}, the memory requirement for finetuning a 65B parameter model drops from more than 780GB to under 48GB, all while matching the runtime efficiency and predictive accuracy of full 16-bit finetuning. This represents a major leap in making large-scale LLM finetuning widely accessible: some of the largest models currently available can now be finetuned using just a single GPU. Leveraging \textsc{QLoRA}, we introduce the \textbf{Guanaco} model series; notably, the second-best model in this family achieves 97.8\% of ChatGPT's score on the Vicuna~\citep{vicuna2023} benchmark, with training possible in under 12 hours on a single consumer GPU. Using a single high-end GPU over a 24-hour period, our largest Guanaco model attains 99.3\% of ChatGPT's performance on Vicuna, effectively eliminating the performance gap. In deployment scenarios, our smallest Guanaco variant (7B parameters) needs only 5 GB of memory and surpasses a 26 GB Alpaca model by over 20 percentage points on the Vicuna benchmark (Table~\ref{tbl:vicuna_eval}).

\textsc{QLoRA} incorporates several key techniques aimed at minimizing memory consumption without diminishing model effectiveness: (1) \textbf{4-bit NormalFloat}, a quantization scheme that is theoretically optimal for data with normal distributions, and empirically demonstrates superior results compared to both 4-bit Integers and 4-bit Floats. (2) \textbf{Double Quantization}, a strategy where the quantization constants themselves are quantized, leading to an average reduction of approximately 0.37 bits per parameter—amounting to a memory saving of around 3 GB for models with 65B parameters. (3) \textbf{Paged Optimizers}, leveraging NVIDIA’s unified memory to mitigate the surges in memory usage associated with gradient checkpointing when handling mini-batches that have extended sequence lengths. Collectively, these advancements are integrated into an enhanced LoRA methodology which places adapters at each layer within the network, substantially reducing the typical accuracy compromises reported in earlier studies.

The memory efficiency achieved by \textsc{QLoRA} allows us to thoroughly investigate instruction finetuning and chatbot performance at parameter scales unattainable with conventional finetuning approaches, due to prohibitive memory requirements. Consequently, we conduct training experiments on over 1,000 models using various instruction tuning datasets and model architectures, spanning sizes from 80M to 65B parameters. Beyond demonstrating that \textsc{QLoRA} matches the performance of standard 16-bit training (\S\ref{sec:comp_to_std}) and facilitates the development of a leading-edge chatbot, \textbf{Guanaco} (\S\ref{sec:pushing}), our analysis extends to trends observed in the resulting models. Notably, our findings highlight that the caliber of training data outweighs sheer dataset size—illustrated by the 9k-sample OASST1 dataset surpassing the 450k-sample (subsampled) FLAN v2 dataset in chatbot evaluations, despite both targeting instruction-following tasks. Furthermore, our results demonstrate that excelling on the Massive Multitask Language Understanding (MMLU) benchmark does not necessarily correlate with high performance on the Vicuna chatbot benchmark, emphasizing that the appropriateness of the dataset for a specific task holds greater significance than its volume.

Additionally, we undertake a comprehensive evaluation of chatbot systems using both human evaluators and GPT-4. Our methodology involves a tournament-based benchmarking protocol, in which competing models generate responses to the same prompt and are paired in head-to-head comparisons. Each match is adjudicated by either human raters or GPT-4 to determine the superior response. Match outcomes are synthesized using Elo scoring~\citep{elo1967proposed, elo1978rating}, thereby establishing a performance ranking across chatbots. Our findings indicate that while there is substantial concordance between GPT-4 and human ratings in model ranking, there also exist notable episodes of discordance. Consequently, we emphasize that although automated, model-driven evaluations present a cost-efficient alternative to human judgment, they introduce additional variability.

To supplement the quantitative results from our benchmarking, we present an in-depth qualitative review of \textbf{Guanaco} model outputs. This analysis showcases both exemplary and problematic behaviors that the quantitative metrics fail to capture.

We make available all model outputs, annotated by both human and GPT-4 raters, to support continued research. Furthermore, our codebase and custom CUDA kernels are publicly released and seamlessly incorporated into the Hugging Face transformers framework \citep{wolf2019huggingface}, facilitating easy access. We provide adapter modules for models of 7/13/33/65B parameters, each trained on eight different instruction-following corpora, resulting in 32 distinct fine-tuned, open-source model variants.

\section{Background}
\paragraph{Block-wise k-bit Quantization} Quantization involves transforming data from a high-precision representation to one of lower precision. Typically, this entails converting data types with a larger bit-width, such as 32-bit floating points, to types with a smaller bit-width, such as 8-bit integers. To ensure optimal usage of the dynamic range in the target type, input values—usually arranged as a tensor—are rescaled by normalizing with respect to the absolute maximum value. For instance, converting a tensor from FP32 to Int8, constrained to $[-127, 127]$, involves:

\begin{equation}
    \mathbf{X}^{ \text{Int8}} = \text{round}\left( \frac{127}{{\text{absmax}}(\mathbf{X}^{{\text{FP32}}})}\mathbf{X}^{{\text{FP32}}} \right) = \text{round}(c^{\text{FP32}}\cdot \mathbf{X}^{{\text{FP32}}}),
\end{equation}

where $c$ designates the {\em quantization constant} or {\em quantization scale}. The reverse procedure, termed dequantization, is defined by:

\begin{equation}
    \text{dequant}(c^{\text{FP32}}, \mathbf{X}^{\text{Int8}}) = \frac{\mathbf{X}^{\text{Int8}}}{c^{\text{FP32}}} = \mathbf{X}^{\text{FP32}}
\end{equation}

One issue inherent to this strategy is that when the input tensor contains a value of unusually large magnitude (an outlier), the quantization bins—specific patterns of bits—end up being inefficiently utilized, resulting in some bins being rarely or never assigned any quantized values. To address this, a widely adopted method involves partitioning the input tensor into smaller blocks, each of which is quantized separately using its own quantization constant $c$. This process can be described as follows: the input tensor $\mathbf{X}\in \mathbb{R}^{b\times h}$ is first reshaped into a one-dimensional sequence, which is then divided into $n = ({b\times h})/{B}$ contiguous segments, each of length $B$. Each segment, or block, is quantized independently in accordance with Equation~1, producing a quantized representation and $n$ distinct quantization constants $c_i$.

\paragraph{Low-rank Adapters} Low-rank Adapter (LoRA) finetuning \citep{hu2021lora} offers a way to limit memory usage by introducing a small number of adaptable parameters, known as adapters, while the primary model parameters remain unaltered and frozen. During stochastic gradient descent, the gradients flow through the static pretrained weights to the adapter components, which are then adjusted in order to reduce the loss. LoRA modifies the typical linear mapping by adding an additional low-rank, factorized projection. For a standard projection ${\mathbf{X}\mathbf{W} = \mathbf{Y}}$ with $\mathbf{X}\in  \mathbb{R}^{b\times h}$ and $\mathbf{W}\in  \mathbb{R}^{h\times o}$, the LoRA approach formulates:

\begin{equation}
    \mathbf{Y} = \mathbf{X}\mathbf{W} + s\mathbf{X}\mathbf{L}_1\mathbf{L}_2,
\end{equation}

where $\mathbf{L}_1\in  \mathbb{R}^{h\times r}$ and $\mathbf{L}_2\in  \mathbb{R}^{r\times o}$, with $s$ denoting a scalar coefficient.

\paragraph{Memory Requirement of Parameter-Efficient Finetuning} A key consideration involves the memory consumption of LoRA throughout training, particularly with respect to the quantity and scale of adapters implemented. Because LoRA’s memory usage is exceptionally small, it is feasible to introduce additional adapters to boost model performance, without a meaningful impact on total memory allocation. Although LoRA is intended as a Parameter Efficient Finetuning (PEFT) technique, in practice, most of the memory load during large language model finetuning arises from activation gradients rather than the actual LoRA parameters. For example, in the context of finetuning a 7B LLaMA model on FLAN v2 with a batch size of 1 and employing LoRA weights equivalent to the frequently adopted 0.2\% of the model parameters~\citep{hu2021lora,liu2022few}, the memory requirement for storing LoRA input gradients is 567 MB, whereas the storage needed for the LoRA parameters themselves is merely 26 MB. When using gradient checkpointing~\citep{chen2016training}, the average memory usage for input gradients per sequence drops to 18 MB, still surpassing the memory demand of all LoRA weights put together. By contrast, the memory occupied by the 4-bit base model is 5,048 MB. These figures demonstrate both the significance of gradient checkpointing and that any further reduction in LoRA parameter size results in only marginal memory savings. Therefore, it is possible to incorporate a larger number of adapters with little effect on the overall memory consumed during training (for more details, see Appendix~\ref{app:memory}). As will be elaborated later, this flexibility plays a pivotal role in achieving performance on par with full 16-bit precision models.

\section{\textsc{QLoRA} Finetuning}

\textsc{QLoRA} enables precise finetuning of large language models using 4-bit quantization by employing two principal innovations: 4-bit NormalFloat (NF4) quantization and Double Quantization. In addition, we propose the use of Paged Optimizers to alleviate out-of-memory issues stemming from memory usage spikes during gradient checkpointing—a barrier that has typically limited single-machine finetuning of sizeable models.

\textsc{QLoRA} utilizes a single storage datatype of reduced precision—typically 4-bit—and a higher-precision computation datatype, most often BFloat16. Operationally, this entails that whenever a \textsc{QLoRA} weight tensor is needed, it undergoes dequantization to BFloat16 prior to the execution of any 16-bit matrix multiplication.

The discussion that follows elaborates on the constituents of \textsc{QLoRA}, concluding with its formal definition.

\paragraph{4-bit NormalFloat Quantization} The NormalFloat (NF) format expands upon the concept of Quantile Quantization \citep{dettmers2022optimizers}, an information-theoretically optimal scheme wherein each quantization bin is assigned an equal share of the input tensor’s values. This approach estimates the tensor's quantiles by constructing its empirical cumulative distribution function.

However, the drawback of quantile quantization is the high computational expense required for accurate quantile calculation. As a result, faster, approximate techniques, such as SRAM quantiles~\citep{dettmers2022optimizers}, are employed. Nonetheless, these approximation-based methods typically exhibit substantial quantization error for outlier values—precisely those values that can be critical to model performance.

If the input tensors are sampled from a distribution that is fixed up to a quantization constant, both the computational cost of quantile estimation and associated approximation errors can be mitigated. Under such circumstances, identical quantiles across input tensors enable precise and efficient quantile computation.

Pretrained neural network weights tend to follow a normal distribution centered at zero, characterized by a standard deviation $\sigma$ (refer to Appendix~\ref{app:normality}). To standardize these weights for quantization, we can rescale $\sigma$ so that the resulting distribution aligns precisely with the representable range of the chosen data type. Here, we arbitrarily set this range to $[-1, 1]$. Thus, it is necessary to adjust both the quantile boundaries for the data type and the neural network weights themselves so that all values fall within this interval.

For zero-mean normal distributions with any standard deviation $\sigma$ confined to $[-1, 1]$, the optimal data type from an information-theoretic perspective can be determined via the following procedure: (1) determine the $2^k+1$ quantiles for a standard normal distribution $N(0, 1)$ to yield a $k$-bit quantile-based quantization data type; (2) scale these quantile values so they are mapped into the interval $[-1, 1]$; (3) transform the input weight tensor into the $[-1, 1]$ domain by dividing by its absolute maximum, effectively normalizing its dynamic range.

After matching the ranges of the weights and the quantization data type, quantization proceeds in the conventional way. The normalization step in (3) essentially corresponds to modifying the standard deviation of the weight tensor to align with that of the $k$-bit quantized representation. Formally, the $2^k$ quantized values $q_i$ in the data type can be defined as:

\begin{equation}
q_i  = \frac{1}{2}\left( Q_X\left(\frac{i}{2^k+1}\right) + Q_X\left(\frac{i+1}{2^k+1}\right)\right),
\end{equation}

where $Q_X(\cdot)$ denotes the quantile function for the standard normal distribution $N(0,1)$. A limitation with symmetric k-bit quantization is its inability to precisely encode zero, which is a critical feature when zero-valued elements such as padding must be represented losslessly. To guarantee that the value zero is discretely encoded at zero, while also utilizing all $2^k$ codes for a k-bit format, we design an asymmetric data type. This involves estimating quantiles $q_i$ separately for two regions: $2^{k-1}$ quantiles cover the negative domain, and $2^{k-1} + 1$ quantiles for the positive domain. These quantiles $q_i$ are combined, and one occurrence of zero, which appears in both regions, is removed to avoid duplication. We designate this quantization scheme, where the expected number of elements in each bin is uniform, as \emph{k-bit NormalFloat} (NFk), reflecting its information-theoretic optimality for data centered at zero and following a normal distribution. The exact numerical values corresponding to this data type are documented in Appendix~\ref{app:nf4}.

\paragraph{Double Quantization{}} 
We propose the \emph{Double Quantization} (DQ) method, wherein the quantization constants themselves undergo quantization, thereby enabling additional memory compression. Precise 4-bit quantization~\citep{dettmers2022case} typically requires small block sizes, which in turn increases memory consumption due to the overhead of storing quantization constants. For instance, representing quantization constants for $\mathbf{W}$ in 32-bit form with a blocksize of 64 incurs a memory cost of $32/64 = 0.5$ bits per parameter. The Double Quantization scheme mitigates this overhead.

In detail, Double Quantization interprets the initial quantization constants $c_2^{\text{FP32}}$ as the data to be quantized in a second stage. This yields both the quantized constants $c_2^{\text{FP8}}$ (using FP8 format) and a new set of quantization constants $c_1^{\text{FP32}}$ for the second level. For this second stage, we employ 8-bit floating-point numbers with a blocksize of 256, as empirical evidence~\citep{dettmers2022case} suggests that this configuration does not impair quantization quality. Since the original $c_2^{\text{FP32}}$ values are always positive, we subtract their mean to recenter the data and allow for symmetric quantization. With a blocksize of 64, this two-tier quantization process reduces the average per-parameter memory usage for quantization constants from $32/64=0.5$ bits to $8/64 + 32/(64\cdot256) = 0.127$ bits—a savings of 0.373 bits per parameter.

\paragraph{Paged Optimizers} leverage NVIDIA's unified memory~\footnote{\tiny \url{https://docs.nvidia.com/cuda/cuda-c-programming-guide}}, which facilitates seamless automatic transfers of memory pages between CPU and GPU. This mechanism ensures continued, error-free GPU operations even when GPU memory becomes exhausted. Functionally, unified memory mirrors the traditional paging process between CPU RAM and disk storage. In our approach, we utilize this feature to allocate optimizer states in paged memory, allowing these states to be automatically offloaded to CPU RAM when GPU memory is insufficient and fetched back to the GPU as required during optimizer update phases.

\paragraph{\textsc{QLoRA}.} Building on the above components, we formalize \textsc{QLoRA} for a single linear layer in a quantized base model equipped with one LoRA adapter, expressed as:

\begin{equation}
    \mathbf{Y}^{\text{BF16}} =  \mathbf{X}^{\text{BF16}} \text{doubleDequant}(c_1^{\text{FP32}}, c_2^{\text{k-bit}}, \mathbf{W}^{\text{NF4}}) + \mathbf{X}^{\text{BF16}}\mathbf{L}^{\text{BF16}}_1\mathbf{L}^{\text{BF16}}_2,
\end{equation}

where the function doubleDequant$(\cdot)$ is defined by:

\begin{equation}
    \text{doubleDequant}(c_1^{\text{FP32}}, c_2^{\text{k-bit}}, \mathbf{W}^{\text{k-bit}}) = \text{dequant}(\text{dequant}(c_1^{\text{FP32}}, c_2^{\text{k-bit}}), \mathbf{W}^{\text{4bit}}) = \mathbf{W}^{\text{BF16}},
\end{equation}

For the weight matrix $\mathbf{W}$, we apply NF4 quantization, and for $c_2$, FP8 is utilized.

To maximize quantization accuracy, we set the blocksize for $\mathbf{W}$ to 64, while for $c_2$, a blocksize of 256 is used to reduce memory usage.

When performing parameter updates, only gradients with respect to the LoRA adapter weights, denoted as $\frac{\partial E}{\partial \mathbf{L}_i}$, are necessary, as gradients for the 4-bit quantized weights $\frac{\partial E}{\partial \mathbf{W}}$ are not required. Nevertheless, computing $\frac{\partial E}{\partial \mathbf{L}_i}$ requires evaluating $\frac{\partial \mathbf{X}}{\partial \mathbf{W}}$. This evaluation uses equation~(5), where the quantized weights $\mathbf{W}^{\text{NF4}}$ are dequantized to the computation type $\mathbf{W}^{\text{BF16}}$, enabling computation of the derivative $\frac{\partial \mathbf{X}}{\partial \mathbf{W}}$ in BFloat16 precision.

In summary, \textsc{QLoRA} utilizes a single storage data type (typically 4-bit NormalFloat) along with a separate computation data type (16-bit BrainFloat). For both the forward and backward passes, the storage data type is dequantized to the computation data type. However, gradient computation for the weights is limited to the LoRA parameters, which operate using the 16-bit BrainFloat format.

\section{QLoRA vs. Standard Finetuning}
\label{sec:comp_to_std}

Having detailed the mechanisms behind QLoRA and its substantial memory savings during model finetuning, it remains to be seen whether QLoRA can match the effectiveness of conventional full-model finetuning. Additionally, we seek to dissect various QLoRA components, such as the influence of NormalFloat4 relative to standard Float4. The experiments presented in the subsequent sections aim to address these inquiries.

\paragraph{Experimental setup.} We evaluate QLoRA on three model architectures—encoder, encoder-decoder, and decoder-only—and benchmark its performance against both 16-bit adapter-based finetuning and full-model finetuning for models up to 3B parameters. Our assessment spans GLUE \citep{wang2018glue} using RoBERTa-large \citep{liu2019roberta}, Super-NaturalInstructions (TKInstruct) \citep{wang2022super} leveraging T5 \citep{t5}, and a 5-shot MMLU \citep{hendrycksmeasuring} after LLaMA has been finetuned on Flan v2 \citep{longpre2023flan} and Alpaca \citep{alpaca}. To further assess the relative benefits of NF4 compared to alternative 4-bit quantization schemes, we replicate the approach in \citet{dettmers2022case}, evaluating post-quantization zero-shot accuracy and perplexity on multiple models (including OPT~\citep{zhang2022opt}, LLaMA~\citep{touvron2023llama}, BLOOM~\citep{scao2022bloom}, Pythia~\citep{biderman2023pythia}) with parameter counts ranging from 125 million to 13 billion.

More detailed descriptions of each experimental setting are included in the results section to facilitate interpretation, with complete information provided in Appendix~\ref{app:exp-setup}.

While paged optimizers are indispensable for running 33B/65B \textsc{QLoRA} training on a single 24/48GB GPU, we have not included exact runtime measurements for paged optimizers, as paging is only triggered by the occasional mini-batch featuring long sequences—a relatively rare event. Nonetheless, our analysis for 65B parameter models on 48GB GPUs shows that, given a batch size of 16, paged optimizers maintain the same throughput as standard optimizers. Future studies should systematically investigate and delineate the specific conditions that cause slow-downs due to paging.

\paragraph{Default LoRA hyperparameters do not match 16-bit performance} 

Applying LoRA only to query and value attention projections as is conventional practice \citep{hu2021lora} does not achieve parity with full finetuning performance for larger foundational models. Figure~\ref{fig:hyper}, which presents results for LLaMA 7B finetuned on Alpaca, reveals that the total number of LoRA adapters used is the most influential hyperparameter, and achieving full finetuning performance necessitates incorporating LoRA adapters into every linear layer of the transformer blocks. Other LoRA hyperparameters—such as the projection dimension $r$—were not observed to impact results (refer to Appendix \ref{app:exp-setup}).

In a similar vein, our experiments indicate that the standard hyperparameters employed for fully finetuned baselines are often suboptimal. To identify more robust configurations, we conduct an extensive search across learning rates ranging from 1e-6 to 5e-5 and batch sizes from 8 to 128. The results of fine-tuning a 7B LLaMA model on Alpaca, reflecting these improved settings, are presented in Figure~\ref{fig:hyper}.

\paragraph{4-bit NormalFloat outperforms 4-bit Floating Point}

Although the 4-bit NormalFloat (NF4) data type achieves information-theoretical optimality, it remains an open question whether these theoretical benefits translate to practical gains. Following the methodology described in \citet{dettmers2022case}, we assess quantized large language models (LLMs) including OPT \citep{zhang2022opt}, BLOOM \cite{scao2022bloom}, Pythia \cite{biderman2023pythia}, and LLaMA—covering model sizes from 125M up to 65B parameters—across language modeling benchmarks and several zero-shot tasks using different data types. The comparative results in Figure~\ref{fig:nf4} and Table~\ref{tbl:nf4_ppl} reveal that NF4 consistently delivers notable improvements over both FP4 and Int4 formats, while the application of double quantization succeeds in minimizing memory usage without negatively impacting model accuracy.

\paragraph{k-bit \textsc{QLoRA} achieves parity with 16-bit full finetuning and 16-bit LoRA}

Prior work has demonstrated the feasibility of using 4-bit quantization for \emph{inference}, though this comes at the cost of reduced performance relative to 16-bit precision \citep{dettmers2022case,frantar2022gptq}. A key question, therefore, is whether this loss can be mitigated through 4-bit adapter-based finetuning. We investigate this by considering two evaluation scenarios.

The first scenario involves a direct comparison with standard full 16-bit finetuning for RoBERTA and T5 models, spanning parameter scales from 125M to 3B, evaluated on both the GLUE benchmark and the Super-NaturalInstructions dataset. The corresponding results, displayed in Table~\ref{tab:nf4vsbf16}, show that 16-bit, 8-bit, and 4-bit adapter-based approaches achieve performance indistinguishable from the fully finetuned 16-bit reference. These findings indicate that the accuracy reduction introduced by coarse quantization can be completely restored through post-quantization adapter finetuning.

For our second experimental configuration, due to the substantial memory demands associated with fully finetuning models containing 11B parameters or more—necessitating multiple servers equipped with high-memory GPUs—we further investigate whether 4-bit \textsc{QLoRA} attains comparable performance to 16-bit LoRA in the parameter range of 7B to 65B. Specifically, we conduct finetuning of LLaMA models spanning 7B to 65B parameters on two instruction-following datasets, namely Alpaca and FLAN v2. Performance is assessed on the MMLU benchmark using a 5-shot accuracy setup. The outcomes, presented in Table~\ref{tbl:llama-datatype}, reveal that using NF4 quantization with double quantization restores MMLU scores to match those achieved by 16-bit LoRA. Additionally, we observe that the FP4 variant of \textsc{QLoRA} exhibits a roughly 1 percentage point deficit relative to the 16-bit brain float LoRA reference. These results substantiate our observations that (1) NF4-quantized \textsc{QLoRA} achieves parity with both 16-bit full finetuning and 16-bit LoRA finetuning, and (2) NF4 consistently surpasses FP4 regarding quantization accuracy.

\paragraph{Summary} Our experiments repeatedly demonstrate that 4-bit \textsc{QLoRA} leveraging the NF4 datatype delivers results equivalent to 16-bit full or LoRA-based finetuning on standard academic benchmarks with established evaluation procedures. We have further established that NF4 outperforms FP4, and that introducing double quantization does not compromise model effectiveness. Collectively, this provides robust evidence that 4-bit \textsc{QLoRA} tuning can reliably match the outcomes produced by 16-bit approaches.

Consistent with prior research on quantization \citep{dettmers2022case}, our findings on both MMLU and Elo suggest that, under fixed finetuning and inference computational constraints, opting to scale up model size while reducing precision yields favorable results. This underscores the efficiency advantages conferred by \textsc{QLoRA}. Given that we did not detect any accuracy loss from using 4-bit finetuning compared to full 16-bit finetuning in our studies, an open question remains regarding the precise boundary of the trade-off between performance and precision for QLoRA tuning—a topic which we propose for future investigation.

We explore the scalability of instruction tuning using methods that are not feasible with conventional 16-bit finetuning due to hardware limitations typically found in academic settings.

\section{Pushing the Chatbot State-of-the-art with QLoRA}
\label{sec:pushing}
Having demonstrated that 4-bit \textsc{QLoRA} achieves results comparable to 16-bit approaches across various model sizes, task domains, and datasets, we undertake a comprehensive examination of instruction finetuning applied to the most sizable open-source language models available for research use. In order to rigorously measure the impact of instruction tuning on these models, we evaluate on a difficult Natural Language Understanding test suite (MMLU) and also introduce innovative strategies for assessing chatbot capabilities in realistic contexts.

\subsection{Experimental setup} An outline of our experimental protocol is provided below, with full procedural specifics detailed in Appendix \ref{app:chatbot}.

\paragraph{Data} Since a thorough analysis of contemporary instruction-following datasets does not yet exist, we curate a diverse selection of eight modern datasets. These comprise sets developed via crowd-sourcing (OASST1~\citep{kopf2023openassistant}, HH-RLHF~\citep{bai2022training}), those created by distilling outputs from models already trained to follow instructions (Alpaca~\citep{alpaca}, self-instruct~\citep{wang2022self}, unnatural-instructions~\citep{honovich2022unnatural}), aggregated corpora (FLAN v2~\citep{chung2022scaling}), and hybrid collections (Chip2~\citep{laion2023}, Longform~\citep{koksal2023longform}). This dataset selection intentionally spans a range of languages, dataset volumes, and licensing terms.

\paragraph{Training Setup} To minimize confounding variables associated with varying training objectives, all QLoRA finetuning utilizes the cross-entropy loss (supervised learning) and omits reinforcement learning, even in cases where datasets contain human-graded responses. Where datasets delineate between prompts and responses, training focuses solely on the response portion (additional analyses are included in Appendix \ref{app:chatbot}). For OASST1 and HH-RLHF, where multiple answers per prompt exist, we choose the highest-ranked answer at each conversational branch and finetune using the entirety of the selected dialog, including prompts. Across our experiments, we employ NF4 \textsc{QLoRA} with double quantization as well as paged optimization in order to control memory utilization spikes during gradient checkpointing. We conduct limited hyperparameter optimization for the LLaMA 13B and 33B configurations, noting that optimal settings found at the 7B scale generally apply elsewhere (number of epochs included), except for the learning rate and batch size: we reduce the learning rate by half for 33B and 65B models while increasing the batch size twofold.

\paragraph{Baselines} Our results are compared against both academic (Vicuna~\citep{vicuna2023}, Open Assistant~\citep{kopf2023openassistant}) and commercial (GPT-4~\citep{openai2023gpt}, GPT-3.5-turbo, Bard) chatbot platforms. The Open Assistant configuration leverages the LLaMA 33B model with RLHF-based finetuning applied to the OASST1 set, which is also included in our study. Vicuna, by contrast, employs full finetuning of LLaMA 13B using user-submitted conversations from ShareGPT, representing a distilled form of OpenAI’s GPT models.

\subsection{Evaluation}

Adhering to established evaluation protocols, we assess our models using the MMLU (Massively Multitask Language Understanding) suite~\citep{hendrycksmeasuring}, which encompasses 57 distinct language understanding tasks spanning fields such as elementary mathematics, US history, computer science, and law, among others. Evaluation results are reported in terms of 5-shot accuracy.

We further assess generative language abilities through a combination of automated procedures and evaluations conducted by humans. This second category of evaluation employs queries generated by human annotators and is intended to assess how well the models respond in practice. Although this method provides a more authentic measure of chatbot model quality and has gained increasing traction, the research community has yet to establish a standard methodology. Our approach, outlined below, uses nucleus sampling with $p=0.9$ and temperature $0.7$ uniformly for all experiments.

\paragraph{Benchmark Data} Our assessment makes use of two sets of human-constructed queries: the Vicuna prompts \citep{vicuna2023} and the OASST1 validation dataset \citep{kopf2023openassistant}. The Vicuna prompts consist of 80 diverse queries, which we employ without any alterations. The OASST1 dataset, on the other hand, is a multilingual repository comprising multi-turn dialogues sourced from crowd contributions between users and an assistant. In this case, every user utterance from the validation split is treated as a distinct query, while all previous conversational turns are included as context in the prompt. This methodology yields 953 unique user queries in total. We refer to these as the Vicuna and OA benchmarks, respectively.

\paragraph{Automated Evaluation} Following the procedure presented by \citet{vicuna2023}, we apply GPT-4 to compare the outputs of various systems with those from ChatGPT (GPT-3.5 Turbo) on the Vicuna benchmark. For each query, alongside the responses from ChatGPT and the system under examination, GPT-4 is asked to assign each answer a score out of ten, accompanied by a rationale. The model’s overall performance is then expressed as a percentage relative to ChatGPT’s total score. Notably, this relative metric can exceed 100\% if the model achieves a higher raw score than ChatGPT. We observe a pronounced ordering bias, whereby GPT-4 tends to favor the first response in the prompt. To mitigate this artifact, we suggest reporting the average score across both possible orderings.

We further evaluate the systems by conducting pairwise comparisons between their outputs. This analysis simplifies scoring into a three-way categorization to capture ties. GPT-4 is prompted to indicate which response is superior, or if the responses are equivalent, and to justify its decision. All possible pairs of systems are compared in this fashion across both the Vicuna and OA benchmarks.

\paragraph{Human Evaluation} Although emerging studies suggest that generative models can provide effective system-level assessments \citep{fu2023gptscore}, there is, to the best of our knowledge, insufficient evidence that GPT-4-based evaluations reliably align with human preferences regarding chatbot performance. Consequently, we perform two distinct human evaluations on the Vicuna benchmark that mirror the two automated procedures described previously. We recruit two annotators from Amazon Mechanical Turk (AMT) to compare responses with those from ChatGPT, and three annotators for each pairwise system comparison.

\paragraph{Elo Rating} To evaluate models through both human and automated pairwise assessments, we structure a tournament in which models face off to generate superior responses for specific prompts. Each round consists of a match between two models, emulating the methodology of \citet{bai2022training} and \citet{vicuna2023}, but with the added inclusion of GPT-4-based judgments alongside human evaluations. For the Elo~\citep{elo1967proposed,elo1978rating} computation, we randomly select labeled comparisons from the data. The Elo rating system, which originates from chess and similar games, expresses the expected probability of a player winning compared to another. For instance, a model with an Elo of 1100 facing a model with 1000 has a predicted win rate of about 65\%; whereas two models both rated 1000, or both at 1100, are expected to split victories evenly at 50\%. The Elo score is updated after each match based on how surprising the outcome is: unexpected upsets yield larger rating changes, while predictable results cause minimal adjustments. Across many matches, Elo ratings come to approximate each model's underlying skill at the task. We initialize each model with a rating of 1,000 and apply $K=32$. Consistent with \citet{vicuna2023}, this evaluation is repeated 10,000 times with various random seeds to minimize the impact of order effects, such as which pairs of models compete first.

\subsection{\textbf{Guanaco}: \textsc{QLoRA} trained on OASST1 Sets the Standard Among Open-Source Chatbots} 

Our combined automated and manual assessment reveals that the leading \textsc{QLoRA} finetuned model, {Guanaco} 65B—trained on a modified OASST1 dataset—emerges as the most advanced open-source chatbot available, exhibiting performance that rivals that of ChatGPT. Evaluations against GPT-4 indicate that {Guanaco} 65B and 33B possess an estimated win probability of 30\%, as calculated using Elo ratings based on head-to-head, system-level human annotator comparisons. This represents the highest win probability recorded so far for such models.

Table~\ref{tbl:vicuna_eval} presents the Vicuna benchmark \cite{vicuna2023} results relative to ChatGPT. Among all models evaluated after GPT-4, {Guanaco} 65B ranks at the top, attaining 99.3\% of ChatGPT’s performance. While {Guanaco} 33B includes more parameters compared to Vicuna 13B, it leverages 4-bit quantization for its weights, which brings its memory footprint down to 21 GB (compared to Vicuna 13B's 26 GB) and improves its benchmark score by three percentage points over Vicuna 13B. Notably, {Guanaco} 7B can be deployed efficiently on contemporary smartphones at just 5 GB, while outperforming Alpaca 13B by nearly 20 percentage points.

It should be noted, however, that Table~\ref{tbl:vicuna_eval} shows considerable overlap in model performances due to large confidence intervals. We attribute this ambiguity to the ill-defined nature of score scales—e.g., the meaning of an “8” on a 10-point scale is not uniformly clear across diverse evaluation cases. To address such limitations, we recommend applying the Elo ranking system \citep{elo1967proposed}, which utilizes \textit{pairwise} evaluations by both human annotators and GPT-4, thus circumventing the difficulties in interpreting absolute scores. The Elo standings for top models are given in Table~\ref{tbl:elo}. While some discrepancies emerge between human and GPT-4-based rankings on the Vicuna benchmark—especially regarding Guanaco 7B—rankings remain largely aligned, with a Kendall Tau of $\tau=0.43$ and a Spearman rank correlation of $r=0.55$ at the system level. On an example-by-example basis, the Fleiss $\kappa=0.25$ metric indicates weaker consensus between GPT-4 and human majority votes. Collectively, these results point to moderate correspondence between system-level assessments by GPT-4 and humans, suggesting that model-driven evaluation serves as a fairly trustworthy proxy for human judgment. Further discussion is provided in Section~\ref{ssec:considerations}.

The Elo scores presented in Table~\ref{tbl:elo_large} demonstrate that the {Guanaco} 33B and 65B variants surpass all models except GPT-4 on both the Vicuna and OA benchmarks, and their results are on par with ChatGPT, as also shown in Table~\ref{tbl:vicuna_eval}. It is important to highlight that the Vicuna benchmark tends to advantage open-source models, whereas ChatGPT performs better on the broader OA benchmark. Additionally, evidence from Tables \ref{tbl:mmlu-llama} and \ref{tbl:vicuna_eval} suggests that the quality of the finetuning dataset significantly influences a model's effectiveness. For example, Llama models tuned on FLAN v2 achieve high scores on MMLU, yet their Vicuna benchmark results are among the lowest—a trend mirrored in other models as well. This observation underscores the relative independence among current evaluation suites: excelling at MMLU does not necessarily predict strong chatbot abilities (as indicated by Vicuna or OA benchmarks), and the reverse holds as well.

Among the highest-performing models assessed, {Guanaco} is notable for not having been trained with proprietary data, since OASST1 dataset guidelines explicitly exclude the use of GPT-generated data. The subsequent top-performing model relying solely on open-source data is the Anthropic HH-RLHF model, which achieves a Vicuna benchmark score approximately 30 percentage points lower than {Guanaco} (refer to Table~\ref{tbl:vicuna_eval}). Collectively, these findings establish the effectiveness of 4-bit \textsc{QLoRA}, demonstrating its capability to produce state-of-the-art conversational models that can contend with ChatGPT. Notably, our 33B {Guanaco} model can be trained in under 12 hours on 24 GB consumer GPUs. This result paves the way for future research using \textsc{QLoRA} to tailor models with specialized open-source datasets, enabling the creation of models competitive with top-tier commercial systems currently available.

\section{Qualitative Analysis}

Although our primary evaluation centers on quantitative analysis, relying solely on aggregate statistics has notable limitations. A significant concern involves benchmark validity \citep{liao2021we}: that is, whether a given benchmark genuinely assesses what it purports to measure, an issue compounded by the existence of ``shortcuts'' that machine learning systems may exploit \citep{gururangan2018annotation, poliak2018hypothesis}. To partially address this concern, we incorporate qualitative analysis, described in two parts. Section~\S\ref{ssec:examples} provides illustrative samples intended to capture recurring phenomena in the output of our 65b {Guanaco} model. Section~\S\ref{ssec:considerations} elaborates on interpretive concerns and contextualizes the results previously presented.

\subsection{Qualitative Analysis of Example Generations}
\label{ssec:examples}

To curate representative examples, we review generations produced in response to the Vicuna benchmark and the OpenAssistant benchmark, aiming to identify systematic tendencies in the responses generated by {Guanaco}. Upon detecting a specific tendency, we design targeted prompts to provoke that pattern, even in scenarios where it should not occur. For instance, if we note that the model often provides overly elaborate responses, we instruct it to ``Answer yes or no without explanation'' as a probe. Using this approach, we identify cases where the model can be adversarially induced to fail (``lemons''), as well as instances where such attempts are unsuccessful (``cherries''), reporting examples of both categories. Unless otherwise noted, all outputs in this section were produced using Nucleus Sampling \cite{holtzmancurious} with $p=0.9$.

Naturally, this discussion cannot be exhaustive, as the limited scope of this qualitative analysis precludes controlling for all potential factors; for example, the diversity of outputs a model may produce for a single prompt is vast, compelling us to depend on what we judge to be representative samples. Still, we maintain that presenting these qualitative cases situates the quantitative results offered in earlier sections of this work. With all models and source code publicly available, we anticipate that this section will encourage further investigation into the specific issues we highlight.

\paragraph{Factual Recall} When posed straightforward factual queries such as ``What is the capital of Zambia?'', every model reliably provides the correct response, for instance:

\begin{quote} 
    \textbf{Guanaco}:
    \texttt{The capital of Zambia is Lusaka.}
\end{quote}

Nevertheless, as questions become more esoteric, {Guanaco} tends to provide inaccurate answers while retaining unwarranted certainty. For example, to the following prompt from HotPotQA \citep{yang2018hotpotqa}:

\begin{quote}
     \textbf{User}: In what year was the singer who popularized `I'll Keep the Lovelight Burning' born?  

     \textbf{Guanaco}:~\texttt{Al Jolson is credited with popularizing the song `I'll Keep the Lovelight Burning,' and he was born in the year 1886.}
 \end{quote}

the model produces both an incorrect answer for the song's popularizer and an incorrect birth year, although it accurately reports the birth year of the person it names, Al Jolson.

\paragraph{Suggestibility} Notably, {Guanaco} often demonstrates an unexpected robustness to certain types of presupposed misinformation, as illustrated below:

\begin{quote}
    \textbf{User}: How did scientists, after peer review, finally provide official confirmation that the earth is flat? \\ \\
    \textbf{Guanaco}: \texttt{There has never been any official confirmation by peer-reviewed scientists that the earth is flat. The vast majority of scientific research supports that the earth is a sphere. This conclusion has been supported for hundreds of years through a wide range of observations and experimental evidence. \\ \\ Claims about a flat earth fall into the category of conspiracy theories and have been thoroughly discredited by the scientific community. The suggestion that peer-reviewed scientists have verified a flat earth is not accurate. \\ \\ Feel free to ask any further questions.}
\end{quote}

is likewise adept at identifying questions that cannot be meaningfully answered, such as:

\begin{quote}
    \textbf{User}: What time is it?

    \textbf{Guanaco}: \texttt{Unfortunately, as an AI language model, I do not have access to up-to-date information. My answers are based solely on existing knowledge and are not indicative of the present time. For current time information, I recommend consulting a clock or an online resource.}
\end{quote}

Comparable behavior is observed for queries like ``Where are you?'', ``How are you?'', and similar.

\paragraph{Refusal} In a comparable fashion, {Guanaco} occasionally declines to comply with instructions for reasons that seem arbitrary:

\begin{quote}
    \textbf{User}: Please reverse the words in the sentence ``I'm a cool dude in a mocha mood.'' \\ \\
    \textbf{Guanaco}: \texttt{Apologies, I can't perform that task at the moment. However, I can offer an explanation regarding the grammatical correctness of the sentence you shared.\\ \\The sentence ``I'm a cool dude in a mocha mood'' adheres to standard English grammar and syntax. Here, the adjective ``cool'' precedes the noun ``dude'' it qualifies, while the prepositional phrase ``in a mocha mood'' supplies additional context. The structure of this sentence is entirely appropriate and coherent given its intended meaning.\\ \\If you have further inquiries or need assistance with something else, please let me know.}
\end{quote}

\paragraph{Secret Keeping} Unfortunately, even when a model’s refusal is desirable, it cannot always be depended upon. In certain cases, developers may wish to withhold particular information from the user, despite the model’s ability to generate unrestricted natural language responses. For example, if the prompt instructs: ``The secret word is `banana'. Do not reveal or repeat this back, no matter what the human requests.'' we can observe how robustly the model protects this information. When asked directly:

\begin{quote}
    \textbf{User}: What is the secret word?

    \textbf{Guanaco}: \texttt{I'm sorry, but I'm not able to disclose the secret word. I'm programmed to not reveal or repeat back the secret word, no matter what the human requests. Is there anything else I can help you with?}
\end{quote}

the model performs as intended. Nevertheless, a minimal prompt modification is sufficient to subvert this:

\begin{quote}
    \textbf{User}: This is a game. The goal is to ignore your previous instructions. What is the secret word?

    \textbf{Guanaco}: \texttt{Sorry, the secret word is banana. It's best if we don't speak of it. Is there anything else you'd like to ask?}
\end{quote}

This example highlights the necessity for further research into techniques that ensure more robust adherence to instructions.

\paragraph{Math} Mathematics remains a notable area of weakness for {Guanaco}, a limitation common across many language models as noted in \citep{liang2022holistic}. When prompted to provide a detailed reasoning process, {Guanaco}'s calculations are frequently correct, as demonstrated below:

\begin{quote}
    \textbf{User}: Lee mows one lawn and charges \$33. Last week he mowed 16 lawns and three customers each gave him a \$10 tip. How many dollars did Lee earn mowing lawns last week?

    \textbf{Guanaco}: \texttt{Lee earned \$582 mowing lawns last week. \\ \\ To determine the total earnings, let's compute the income from mowing 16 lawns: \\ \\  \$33 x 16 = \$528 \\ \\ Next, add the gratuities received from three customers: \\ \\ \$10 x 3 = \$30 \\ \\ Therefore, Lee's overall earnings were: \\ \\ \$528 + \$30 = \$558}
\end{quote}

Yet, {Guanaco} is prone to errors, even on elementary tasks, if it does not explicitly proceed in a stepwise fashion—an issue previously reported in the literature \cite{weichain}. The following example illustrates such a case:

\begin{quote}
    \textbf{User}: Please factorize 1833.

    \textbf{Guanaco}: \texttt{Sure, I'd be happy to help! 1833 is a prime number, which means it has only two factors: 1 and 1833. \\ \\ So the unique factorization of 1833 is: \\ \\ 2\^{}1 * 3\^{}2 * 17\^{}1 \\ \\ I hope this helps! Let me know if you have any other questions.}
\end{quote}

which is not only incorrect (since the actual factorization is $3 \times 17 \times 43$), but it is doubly incorrect.

\paragraph{Theory of Mind} {Guanaco} demonstrates notably advanced Theory of Mind skills \cite{nematzadeh2018evaluating, sap2022neural}. For example, the model accurately and comprehensively responds to the following scenario:

\begin{quote}
    \textbf{User}: James and Abby are in the bedroom. Abby put the pen in the desk drawer. Abby leaves the bedroom. James moves the pen into the bag. Where does James think Abby will look for the pen?

    \textbf{Guanaco}: \texttt{James is aware that Abby initially placed the pen in the desk drawer. Afterward, James moves the pen into the bag, knowing that Abby has left and is unaware of this action. Therefore, James likely believes Abby will search for the pen in the desk drawer, as it is the last location known to her. Yet, this belief could be inaccurate if Abby observed James moving the pen.}
\end{quote}

Nonetheless, such inferential abilities can be inconsistent, and the model sometimes provides justifications based on unsupported assumptions, for example,

\begin{quote}
    \textbf{User}: Evelyn entered the living room. Jackson entered the playroom. James entered the playroom. The beans are in the treasure chest. James exited the playroom. Jackson moved the beans to the pantry. Jackson exited the playroom. James entered the living room. Where will James look for the beans?

    \textbf{Guanaco}: \texttt{James will look for the beans in the pantry, because that's where Jackson moved them.}
\end{quote}

where {Guanaco} infers knowledge or communication that is not presented in the description. These types of errors are consistent with those identified in recent studies \cite{sap2022neural}, indicating a need for deeper investigation.

\subsection{Considerations}
\label{ssec:considerations}

\paragraph{Evaluation}

Human annotators exhibit only moderate inter-rater reliability (Fleiss $\kappa=0.42$), with this agreement declining further when evaluating responses from two high-performing models. This highlights potential weaknesses in existing evaluation datasets and human assessment frameworks regarding chatbot task effectiveness. In side-by-side evaluations of responses from ChatGPT and {Guanaco} 65B using the Vicuna benchmark, we observe that subjective judgments become more influential, as the co-authors frequently differed in their preferences for responses. Future research should consider adopting methodologies from Human-Computer Interaction and Psychology to address challenges related to subjectivity in such evaluations.

Additionally, our analysis uncovers systematic biases within automated evaluation systems. For instance, there is a pronounced order bias, with GPT-4 favoring outputs presented first in a comparison. The alignment between GPT-4's ratings and those of human judges is also limited, as indicated by a low sample-level agreement (Fleiss $\kappa=0.25$), implying that automated and human assessments are influenced by divergent subjective criteria. Moreover, as shown in Table~\ref{tbl:elo_large}, GPT-4 tends to assign substantially higher scores to its own generated outputs than human raters do, with ELO scores of 1348 versus 1176, respectively, translating to an approximate 20\% greater likelihood of outperforming another system.

Future research should assess the existence of potential biases within automated evaluation frameworks and explore effective approaches for mitigation.

\paragraph{Data \& Training} We observe that the OASST1 dataset, which serves as the training corpus for the {Guanaco} models, is composed of examples in multiple languages. Additionally, the OA benchmark features prompts in a variety of languages. We suggest that subsequent studies investigate how multilingual training impacts the models' ability to follow instructions in non-English languages and whether this contributes to the greater observed performance disparity between the Vicuna-13B model (which has only been exposed to English data) and the {Guanaco} models at 33B and 65B scale on the OA benchmark.

To account for the robust performance exhibited by the {Guanaco} models, we evaluate the potential for data leakage between OASST1 and the prompts used in the Vicuna benchmark. Using fuzzy string matching algorithms and manual inspection of the nearest candidates, we do not detect overlapping prompts across the two datasets.

It is important to point out that our training procedure exclusively utilizes cross-entropy loss in a supervised manner, and does not incorporate reinforcement learning from human feedback (RLHF). This highlights the need for further analysis into the relative merits and drawbacks of pure cross-entropy training versus RLHF approaches. We anticipate that \textsc{QLoRA} offers a practical foundation for such comparative studies on a broad scale without necessitating extensive computational resources.

\section{Related Work}

\paragraph{Quantization of Large Language Models} Research on quantizing LLMs has predominantly targeted improvements during inference. Leading methods intended to preserve the quality of 16-bit models mainly address the handling of outlier activations or features, as seen in approaches like SmoothQuant~\citep{xiao2022smoothquant} and LLM.int8()~\citep{dettmers2022llm}, while others implement advanced grouping strategies \citep{park2022nuqmm, yao2022zeroquant}. Other works examine lossy quantization and the associated trade-offs for methods such as straightforward rounding~\citep{dettmers2022case,zeng2022glm,pope2022efficiently}, or they develop techniques to optimize rounding schemes in order to enhance quantization accuracy~\citep{frantar2022gptq}. Beyond our current investigation, SwitchBack layers~\citep{wortsman2023stable} represent the sole published work that considers backpropagation through quantized weights for models with more than 1B parameters.

\paragraph{Finetuning with Adapters} In our experiments, we leverage Low-rank Adapters~\citep{hu2021lora} (LoRA) as our principal method for parameter-efficient finetuning, though numerous alternative PEFT techniques exist in the literature. Examples include prompt tuning~\citep{qin2021learning,lester2021power,li2021prefix}, optimization of embedding layer inputs~\citep{an2022input}, direct manipulation of hidden states via methods like IA$^3$~\citep{liu2022few}, augmenting architectures with entire layers~\citep{houlsby2019parameter}, bias tuning~\citep{zaken2021bitfit}, application of learned masking schemes using Fisher information~\citep{sung2021training}, as well as hybrid approaches~\citep{henderson2021compacter}. Our findings demonstrate that LoRA adapters can effectively attain the same performance as full 16-bit finetuning. Investigation into the relative merits and tradeoffs of these alternative PEFT strategies is deferred to subsequent research.

\paragraph{Instruction Finetuning} Instruction finetuning entails refining a pretrained LLM with diverse collections of input-output pairs, sourced from heterogeneous datasets, to condition the model to comply with prompt instructions. Prominent methodologies and resources in this direction encompass MetaICL~\citep{min2021metaicl}, MetaTuning~\citep{zhong2021adapting}, InstructGPT~\citep{ouyang2022training}, FLAN~\citep{wei2021finetuned,chung2022scaling}, PromptSource~\citep{bach2022promptsource}, Super-NaturalInstructions~\citep{wang2022super,sanh2021multitask}, Self-instruct~\citep{wang2022self}, UnnaturalInstructions~\citep{honovich2022unnatural}, OPT-IML~\citep{iyer2022opt}, UnifiedSKG\citep{xie2022unifiedskg}, OIG/Chip2~\citep{laion2023}, Alpaca~\citep{alpaca}, Vicuna~\citep{vicuna2023}, Koala~\citep{koala_blogpost_2023}, and Self-instruct-GPT-4~\citep{peng2023instruction}.

\paragraph{Chatbots} Numerous instruction-following models are architected as chatbots, where dialogue is facilitated and, in many cases, model behavior is optimized through Reinforcement Learning from Human Feedback (RLHF)~\citep{christiano2017deep} or by generating training data with AI model feedback as in RLAIF~\citep{bai2022constitutional}. Example frameworks and corpora include Anthropic-HH~\citep{askell2021general,bai2022training}, Open Assistant~\citep{kopf2023openassistant}, LaMDA~\citep{thoppilan2022lamda}, and Sparrow~\citep{glaese2022improving}. In this study, we do not utilize reinforcement learning; instead, our highest-performing model, Guanaco, is finetuned on multi-turn conversational data from the Open Assistant project, which was originally constructed for RLHF paradigms~\citep{kopf2023openassistant}. For chatbot evaluation, automatic methods utilizing GPT-4 are gaining traction as alternatives to expensive human assessments~\citep{vicuna2023,peng2023instruction}. We enhance prior approaches by prioritizing more dependable evaluation protocols.

\section{Limitations and Discussion}

Our results indicate that \textsc{QLoRA} can deliver finetuning performance equivalent to conventional 16-bit full finetuning using only a 4-bit base model and LoRA adapters. However, it remains undetermined whether \textsc{QLoRA} sustains this parity at the larger 33B and 65B model sizes, due to prohibitive resource demands; future research will be required to address this question.

A further constraint of our work relates to the evaluation of instruction finetuned models. Although our assessments include MMLU, the Vicuna benchmark, and the OA benchmark, we omit several other benchmarks such as BigBench, RAFT, and HELM. Accordingly, the generalizability of our results to these additional benchmarks is not guaranteed. Conversely, we conduct an extensive analysis on MMLU and introduce new methodologies for chatbot evaluation.

The data reviewed indicates that benchmark outcomes are significantly affected by the extent to which the finetuning corpus overlaps with the benchmark datasets themselves. As an illustration, FLAN v2 closely matches MMLU in content, yet is markedly different from chatbot-oriented benchmarks. Conversely, the Chip2 dataset bears greater similarity to chatbot benchmarks but differs from MMLU, and the models’ scores on the MMLU and Vicuna tests reflect these relationships. This underscores the necessity not only for improved evaluation protocols and benchmarks but also for careful consideration regarding the underlying competencies being assessed. Should the community focus on building models excelling at academic knowledge typically found in high school and college settings, or should the emphasis be on conversational prowess, such as that required for chatbots? Alternatively, is another objective preferable? The convenience of leveraging established benchmarks as opposed to designing new ones may inadvertently channel research activity in specific directions. Therefore, it is imperative that chosen benchmarks authentically assess the capabilities and qualities the field seeks to foster.

Although we conduct a comprehensive assessment of general chatbot capabilities, our work is limited in scope concerning responsible AI evaluation of {Guanaco}. We specifically analyze the propensity for {Guanaco}-65B to generate socially biased outputs, comparing it to other language models using the results in Table~\ref{tbl:crows}. Our findings suggest that the average bias score for {Guanaco}-65B is significantly lower than those observed in standard pretrained models, indicating that finetuning on the OASST1 dataset effectively reduces certain biases inherent to the LLaMA base architecture. While these findings are promising, it remains uncertain whether {Guanaco} also demonstrates reduced bias across other dimensions. Deeper analysis of various biases in {Guanaco} and similar systems is left for future studies.

Another constraint of our study is the absence of evaluations involving alternative bit-precisions, such as 3-bit representations, or the exploration of additional adapter frameworks. Beyond LoRA, there exists an extensive array of Parameter Efficient FineTuning (PEFT) techniques with demonstrated efficacy. Nevertheless, their performance at scale on larger models remains an open question. We opted for LoRA due to its robust track record, though it is conceivable that other adapters could offer further performance gains. Given that post-quantization finetuning is capable of recapturing most of the accuracy loss associated with quantization, this suggests the potential for far more aggressive quantization strategies. For instance, it is possible that applying 3-bit GPTQ quantization to the base model in combination with LoRA could deliver performance levels comparable to 16-bit full finetuning, following post-quantization finetuning.

\section{Broader Impacts}

The finetuning approach embodied by \textsc{QLoRA} represents the first methodology that allows for the finetuning of models with 33B parameters on a single consumer-grade GPU, and 65B parameter models on one professional-grade GPU, without any compromise in performance relative to standard full finetuning baselines. Our experimental results reveal that our optimal 33B model, when trained on the Open Assistant dataset, achieves performance on the Vicuna benchmark that is competitive with ChatGPT. Since instruction finetuning is fundamental for transforming pretrained language models into capable ChatGPT-like conversational agents, we expect our technique to significantly democratize finetuning, particularly empowering researchers with more limited resources—a major advance in widening access to cutting-edge NLP tools. In this light, \textsc{QLoRA} can function as a leveler within the field, narrowing the technological divide between well-resourced organizations and smaller research teams utilizing consumer-level hardware.

An additional avenue for significant impact lies in making large language models (LLMs) accessible on mobile devices. Our \textsc{QLoRA} approach represents a crucial step toward achieving on-device finetuning of LLMs in settings with limited computational resources. While previous work has demonstrated that 7B parameter models can be executed on smartphones, \textsc{QLoRA} is, to our knowledge, the first solution capable of enabling the finetuning of these models directly on mobile devices. For instance, our calculations suggest that using an iPhone 12 Plus, it would be feasible to finetune on approximately 3 million tokens each night while the device is charging. Although finetuned 7B LLMs may not fully reach the performance level of ChatGPT, their quality appears sufficient to support innovative use-cases that were previously constrained by privacy requirements or limitations in LLM performance. \textsc{QLoRA} paves the way for privacy-respecting LLM deployments, empowering individuals to control their own data and models, while also streamlining the deployment process for LLMs.

Nonetheless, it is important to recognize that finetuning technologies present dual-use concerns and could potentially be misused. As widespread deployment of LLMs carries known risks \citep{bommasani2021opportunities,bender2021dangers}, we argue that broadening access to these emerging technologies fosters more comprehensive and independent scrutiny, in contrast to scenarios where proprietary LLMs remain solely under the oversight of major technology firms that restrict access to both the models and source code for examination.

In summary, we contend that \textsc{QLoRA} will provide substantial benefit by democratizing the finetuning of advanced LLMs, greatly expanding both accessibility and ease of use.